\subsection{Additional score identities}\label{app:score-details}
Under a shared affine transformation $z\mapsto Az+b$, centering removes
the additive term $b$ but leaves the transformed contrast $Ad$.

The estimator has an exact decomposition. Write
$e_{msak}=d_{msak}-\widehat\beta_{msk}r_a$. Then
\begin{equation}
 \widehat D_{ms}=
 (\widehat\beta_{ms1}-1)(\widehat\beta_{ms2}-1)
 +\frac1H\sum_a\langle e_{msa1},e_{msa2}\rangle.
 \label{eq:decomposition}
\end{equation}
The first term measures reference-aligned amplitude error. The residual is
orthogonal to the single stacked reference configuration; it is not necessarily
outside the feature-space span of the reference means, much less outside all
authentic artistic variation. Both terms can be negative in finite samples.

Let $\theta_s$ stack the four conditional mean contrasts and
$\bar\theta=S^{-1}\sum_s\theta_s$. The population target decomposes as
\begin{equation}
 D_{\rm scene}=\frac{\|\bar\theta-r\|^2}{H}
 +\frac1{SH}\sum_s\|\theta_s-\bar\theta\|^2
 =D_{\rm aggregate}+V_{\rm scene}.
 \label{eq:scene-decomposition}
\end{equation}
To separate contrast strength from direction, define
\begin{equation}
 \widehat Q_m=\frac1{SH}\sum_{s,a}\langle d_{msa1},d_{msa2}\rangle,
 \qquad \widehat D_m=1-2\widehat\beta_m+\widehat Q_m.
 \label{eq:magnitude}
\end{equation}
For positive $\widehat Q_m$, $\widehat\beta_m/\sqrt{\widehat Q_m}$ is a
noise-corrected alignment summary; finite estimates need not lie in $[-1,1]$.
Scaling every centered generated contrast by $c$ gives
$\widehat D_m(c)=1-2c\widehat\beta_m+c^2\widehat Q_m$.
We fit $c=\max(0,\widehat\beta/\widehat Q)$ on the other 13 scenes and evaluate
it on the held-out scene. The fitting criterion is available only when its
training $\widehat Q$ is positive, as it is in every observed fold. This
counterfactual concerns feature vectors, not an implemented image transformation.

\subsection{Paired comparisons and declared sensitivities}\label{app:declared-sensitivities}
We fit the scene fixed-effects regression
\begin{equation}
 \widehat D_{ms}=\alpha_s+\gamma_m+\varepsilon_{ms},
 \qquad\gamma_{\text{GPT Image 2}}=0.
 \label{eq:regression}
\end{equation}
In the balanced panel, coefficient differences equal the paired scene means
used for inference. Changing the regression baseline leaves these differences
unchanged; Table~\ref{tab:specificity-models} presents them relative to FLUX.
Declared sensitivities delete each scene, resample scenes
and reference works, examine feature families and omit texture. Central-square
windows control aspect ratio by discarding peripheral content in generated and
reference images. FLUX retains the lowest uncalibrated error after every scene
deletion ($D=.754$--$.852$) and with square windows ($D=.765$; next-lowest
Nano Banana 2, $1.024$).

\subsection{Distributional proximity and repeat spread}
For painter $a$, energy discrepancy compares its generated collection $Y_a$
with reference collection $X_a$:
\begin{equation}
 \E(X,Y)=\frac{2}{n_Xn_Y}\sum_{i,j}\|x_i-y_j\|
 -\frac1{n_X^2}\sum_{i,i'}\|x_i-x_{i'}\|
 -\frac1{n_Y^2}\sum_{j,j'}\|y_j-y_{j'}\|.
 \label{eq:energy}
\end{equation}
Its finite-sample expectation depends on within-collection spread; equal
generated sample sizes do not eliminate comparison bias. Appendix~\ref{app:diagnostics}
gives the correction for the repeated-scene design. Within-scene repeat spread,
$\frac{1}{2S}\sum_s\|z_{msa1}-z_{msa2}\|^2$, is the mean sample-variance
trace of the two outputs. It separates repeat dispersion from variation across
briefs.
